\documentclass[12pt]{exam}

% --- Packages ---
\usepackage[utf8]{inputenc}
\usepackage[T1]{fontenc}
\usepackage[french]{babel}
\usepackage{amsmath, amssymb}
\usepackage{geometry}
\usepackage{tikz}
\usepackage[most]{tcolorbox}
\usepackage{lmodern}
\usepackage{pgfplots}
\pgfplotsset{compat=1.18}
\usepackage{mdframed}
\usepackage{bm}
\usepackage{xcolor}
\usepackage{sectsty}
\usetikzlibrary{calc}

\usepackage[dvipsnames]{xcolor}
\usepackage{tkz-tab}
\usepackage{multido}

\geometry{margin=2cm}
\setlength{\parindent}{0pt}

% =====================
% Document
% =====================
\begin{document}

\begin{center}
    {\Large \textbf{Exercices demi-groupe}}\\[2mm]
    {\large 2nde}\\[2mm]
    \textit{Équations et inéquations}
\end{center}

\vspace{0.3cm}

\textbf{Consigne générale :}  
Résoudre les équations et inéquations suivantes. Les résultats seront donnés sous forme d’ensembles de solutions.

\vspace{0.4cm}

% ==================================================
\begin{questions}

% -----------------------------
\question \textbf{Équations simples}

Résoudre les équations suivantes :
\begin{parts}
\part $3x - 5 = 7$
\vspace{0.3cm}
\part $2(x - 4) = 5x + 1$
\vspace{0.3cm}
\part $\dfrac{4x - 1}{3} = 2$
\end{parts}

\vspace{0.5cm}

% -----------------------------
\question \textbf{Inéquations simples}

Résoudre les inéquations suivantes :
\begin{parts}
\part $5x + 3 \leq 2x - 6$
\vspace{0.3cm}
\part $-3(x - 2) > 9$
\end{parts}

\vspace{0.5cm}

% -----------------------------
\question \textbf{Équations produits et quotient}

Résoudre les équations suivantes :
\begin{parts}
\part $(2x - 3)(x + 1) = 0$
\vspace{0.3cm}
\part $\dfrac{x-4}{3x+2} = 0$
\end{parts}

\vspace{0.5cm}

% -----------------------------
\question \textbf{Inéquations produits}

Résoudre les inéquations suivantes à l’aide d’un tableau de signes :
\begin{parts}
\part $(x - 2)(x + 3) \geq 0$
\vspace{0.3cm}
\part $(2x + 1)(x - 4) < 0$
\end{parts}

\vspace{0.5cm}

% -----------------------------
\question \textbf{Inéquations quotients}

Résoudre les inéquations suivantes à l’aide d’un tableau de signes :
\begin{parts}
\part $\dfrac{x - 1}{x + 2} \geq 0$
\vspace{0.3cm}
\part $\dfrac{2x + 3}{x - 4} < 0$
\end{parts}

\vspace{0.5cm}

% -----------------------------
\question \textbf{Pour aller plus loin}

Résoudre l’inéquation suivante :
\[
\dfrac{(x - 2)(x + 1)}{x - 3} \geq 0
\]

% -----------------------------
\question \textbf{Vers la 1ère}

Les questions sont indépendantes.

\begin{parts}
    \part Expliquer pourquoi si $x \geq 3$, alors $3x-4$ est positif.
    \part Expliquer pourquoi si$t>-6$ alors $t^2+t+7>0$.
    \part En partant de $1<2$, expliquer pourquoi $2^n < 2^{n+1}$ pour 
    tout $n \in \mathbb{N}$.
    \part On considère l'inéquation $n^3 + n+1 \geq 100$ où l'inconnue 
    $n$ est un entier naturel.
    \\
    Chercher avec la calculatrice la plus petite valeur de $n$ solution de l'inéquation.
\end{parts}

\end{questions}

\end{document}
